
\chapter{Introduction}
\label{ch:intro}

Machine-learning models have reached state-of-the-art skill in weather
forecasting, rivaling traditional physics-based numerical models at a fraction
of their computational cost
\parencite{lam2023graphcast,bi2023pangu}.
More recently, generative models have extended these advances to probabilistic
forecasting
\parencite{price2023gencast,couairon2024archesweather}.
Instead of producing a single deterministic forecast, they are trained to
approximate the conditional distribution of future atmospheric states given
the available atmospheric state and can therefore generate multiple plausible
future trajectories
\parencite{price2023gencast}.
This allows them to represent forecast uncertainty explicitly, alleviating the
smoothing effects often observed in deterministic neural forecasts, enabling
direct ensemble generation, and achieving strong probabilistic skill as
measured, for instance, by the Continuous Ranked Probability Score (CRPS)
\parencite{hersbach2000crps,price2023gencast}.

Beyond their use as probabilistic samplers, generative models can also be
actively steered during generation.
In image generation, applications such as text-to-image generation, editing,
and inpainting routinely rely on \emph{guidance} to steer a pretrained
generative model toward a desired outcome
\parencite{dhariwal2021diffusion,ho2022classifierfree,nichol2021glide}.
Generative weather models, by contrast, are still predominantly used in their
unguided form, while their potential for controlled generation remains
comparatively unexplored.
Only a small number of recent studies have investigated related forms of
weather control, including classifier-guided generation of tropical cyclones
with Climate in a Bottle \parencite{brenowitz2025climate},
precipitation-reduction interventions through guided diffusion
\parencite{ueyama2026guided}, and optimized heatwave storylines using
differentiable weather models
\parencite{whittaker2025heatwaves}.

This opens up a complementary use of generative weather models beyond
probabilistic forecasting.
While forecasting asks which future weather trajectories are likely given the
current state, \emph{scenario generation} asks what a prescribed alternative
evolution could look like while remaining meteorologically plausible
\parencite{shepherd2018storylines}.
This distinction is particularly relevant for extreme weather.
Extreme events account for a disproportionate share of weather-related impacts
\parencite{wmo2023atlas}, yet their rarity means that both historical
observations and finite forecast ensembles provide only limited examples from
which to study them.
The historical record necessarily contains only a finite sample of possible
extremes, while risk analysis may also require considering plausible
alternatives, including scenarios more severe than those previously observed
\parencite{thompson2017unseen,shepherd2018storylines}.

We are therefore interested in exploiting the ability of generative weather
models to produce coherent weather trajectories while deliberately steering
them toward prescribed extremes.
Rather than generating an isolated extreme state, the objective is to obtain an
event that remains structured as weather throughout its evolution: coherent
over time and space, across atmospheric pressure levels, and between the
different meteorological variables.
In other words, the imposed extreme should remain embedded within a
multivariate atmospheric evolution that is consistent with the structure of
realistic weather.

Transferring guidance methods from image generation to this setting is not
immediate.
First, the conditioning objective is different.
Image-generation guidance commonly targets semantic properties, whereas a
weather scenario can be specified quantitatively by prescribing the evolution
of selected atmospheric quantities over a spatial domain and forecast horizon
\parencite{dhariwal2021diffusion,ho2022classifierfree,nichol2021glide}.
Second, plausibility is considerably more structured.
Satisfying such a target is not sufficient on its own: the generated trajectory
must remain coherent beyond the prescribed quantities and spatial domain,
across neighboring locations, atmospheric variables, pressure levels, and
subsequent forecast steps.
Third, the role of guidance strength changes.
In image generation, guidance strength is commonly used to control the
trade-off between stronger conditioning and sample fidelity or diversity
\parencite{dhariwal2021diffusion,ho2022classifierfree}.
For weather scenario generation, an explicit quantitative target instead gives
us a measurable objective: the model should be guided sufficiently to reach
the prescribed trajectory, while the consequences of that guidance on the
remaining forecast must be evaluated explicitly.

These considerations lead to the central research question of this thesis:

\begin{quote}
\emph{Can we steer generative probabilistic weather models to produce extreme
weather trajectories that satisfy a prescribed target while remaining
meteorologically plausible?}
\end{quote}

We address this question through a general framework composed of three stages:
\emph{Specify}, \emph{Guide}, and \emph{Evaluate}.
The first determines the target trajectory toward which the model should be
steered.
The second modifies the generative process to reach that target.
The third assesses both whether the prescribed event has been realized and
whether the resulting trajectory remains consistent with the structure of
realistic weather.

We investigate this framework using ArchesWeather
\parencite{couairon2024archesweather}, a probabilistic weather model based on
flow matching.
As an experimental case study, we focus on regional surface-temperature
heatwaves and formulate the guidance objective in terms of a prescribed
temperature trajectory over a chosen region.
These choices define the particular scenario studied in this thesis rather
than the methodology itself: the guidance framework operates on differentiable
objectives defined on the forecast state and can therefore, in principle, be
used to prescribe different meteorological quantities and forms of weather
evolution.

The three stages of the framework correspond to the main contributions of this
thesis:

\begin{enumerate}[noitemsep]

\item \textbf{Specify.}
We formalize the desired scenario as a quantitative target trajectory of a
selected atmospheric quantity over a prescribed spatial region.
We consider two complementary constructions.
For controlled evaluation of the guidance method, targets are defined as
prescribed relative deviations from a matched unguided forecast.
For storyline generation, targets are constructed from data-driven historical
weather trajectories, providing complete temporal evolutions toward which the
model can be steered.

\item \textbf{Guide.}
We develop a target-driven, gradient-based guidance method for
flow-matching weather models.
Rather than selecting a fixed guidance strength and accepting the resulting
level of extremeness, guidance is formulated around an explicit target
trajectory and adapted to drive the remaining gap toward zero.
Guidance strength is therefore treated as the means of reaching a specified
scenario rather than as the scenario definition itself.

\item \textbf{Evaluate.}
We develop a set of diagnostics for both guidance behavior and weather
plausibility.
Beyond measuring whether the target trajectory is reached, we evaluate how
the generated scenarios behave across space, pressure levels, variables,
and forecast time.
In particular, we construct a historical benchmark in which observed
weather states provide plausibility anchors and the corresponding unguided
forecast provides a baseline for quantifying the effects of guidance.

\end{enumerate}

If successful, such a framework enables several forms of controlled extreme
scenario generation.
Starting from a given weather state, one may generate alternative future
trajectories at different prescribed levels of extremeness.
Data-driven extreme trajectories may instead be used as explicit targets for
constructing weather storylines conditioned on a particular initial state.
In either case, multiple trajectories and intensity levels can be sampled to
explore not a single deterministic extreme, but a family of guided scenarios.
Although extreme events provide the focus of this thesis, the underlying idea
is more general: guidance can be used to explore alternative weather
trajectories that deliberately depart from the base forecast toward any
suitably defined target.

The remainder of the thesis is organized as follows.
\Cref{ch:background} introduces the technical background on flow matching and
guidance, reviews existing approaches to controlled weather generation, and
describes ArchesWeather.
\Cref{ch:specify,ch:guide,ch:eval} develop the three stages of the framework in
turn: specifying the target trajectory, designing the guidance method, and
defining the evaluation framework.
\Cref{ch:experiments} presents the experimental results, which are discussed in
\Cref{ch:discussion}.
Finally, \Cref{ch:conclusion} summarizes the main findings and outlines
directions for future work.
